% © 2026 Ing. David Jaroš — CC BY-NC-ND 4.0
%
% This work is licensed under a Creative Commons Attribution-NonCommercial-NoDerivatives
% 4.0 International License (CC BY-NC-ND 4.0).
%
% License History: Earlier drafts (up to v0.3) were released under CC BY 4.0.
% From v0.4 onward, all material is released under CC BY-NC-ND 4.0 to protect
% the integrity of the theoretical work during ongoing academic development.

\documentclass[12pt,a4paper]{article}
\usepackage{amsmath,amssymb,amsthm}
\usepackage{mathtools}
\usepackage{geometry}
\usepackage{booktabs}
\usepackage{hyperref}
\usepackage{xcolor}
\usepackage{tcolorbox}
\tcbuselibrary{skins,breakable}
\geometry{left=2.5cm,right=2.5cm,top=2.5cm,bottom=2.5cm}

\newtheorem{theorem}{Theorem}[section]
\newtheorem{lemma}[theorem]{Lemma}
\newtheorem{proposition}[theorem]{Proposition}
\theoremstyle{remark}
\newtheorem{remark}[theorem]{Remark}

\newcommand{\Neff}{N_{\mathrm{eff}}}
\newcommand{\Nhelicity}{N_{\mathrm{helicity}}}
\newcommand{\Ncharge}{N_{\mathrm{charge}}}

\newtcolorbox{statusbox}[1][]{colback=gray!8!white,colframe=black!60,
  arc=3pt,boxrule=1pt,fonttitle=\bfseries,title=Status Summary,#1}

\title{Derivation Audit of $N_{\mathrm{helicity}} = 2$\\
       \large Three Paths and Their Outcomes from $S_{\mathrm{kin}}[\Theta]$}
\author{Ing.\ David Jaro\v{s}}
\date{2026-05-10}

\begin{document}
\maketitle

\begin{abstract}
The factor $\Nhelicity = 2$ appears in Route~R1 and Route~R4 of
\texttt{neff12\_derivations.tex} as part of the $N_{\mathrm{eff}} = 12$
decomposition.  This document examines three candidate derivations
of this factor from $S_{\mathrm{kin}}[\Theta]$ and determines whether
any of them constitutes a genuine loop-count multiplier.
\textbf{Verdict}: none of the three paths derives $\Nhelicity = 2$ as an
additional independent multiplier from the minimal scalar action.
The loop-safe value is $\Nhelicity = 1$ from $S_{\mathrm{kin}}[\Theta]$,
giving $\Neff(\mathrm{loop}) = 3$ and $B_0 = 2\pi$.
\end{abstract}

\tableofcontents

%──────────────────────────────────────────────────────────────────────────────
\section{Setup and Conventions}
\label{sec:setup}
%──────────────────────────────────────────────────────────────────────────────

The audited kinetic action is (\texttt{canonical/n\_eff/step2\_AUDIT.tex}, §A):
\begin{equation}
  S_{\mathrm{kin}}[\Theta]
  = \int d^4x\,d\psi\;
  \mathrm{Sc}\!\bigl[(D_\mu\Theta)^\dagger(D^\mu\Theta)\bigr],
  \label{eq:S_kin}
\end{equation}
which reduces after $\psi$-integration to
\begin{equation}
  S_{\mathrm{eff}} = \sum_{k=1}^{3}
  \int d^4x\;\bigl|({\partial}_\mu + ieA_\mu)\Phi_k\bigr|^2,
  \qquad \Phi_k = \Theta_k \in \mathbb{C}.
  \label{eq:S_eff}
\end{equation}
This is \emph{three copies of scalar QED} for a charged complex scalar.

The reference coefficient for one complex scalar (Peskin \& Schroeder,
scalar QED, Eq.~(7.90) adapted):
\begin{equation}
  B_0^{(1\text{ complex scalar})} = \frac{2\pi}{3}.
  \label{eq:baseline}
\end{equation}
In the convention of \eqref{eq:S_kin}:
$1/\alpha(\mu) = 1/\alpha(\mu_0) + (B_0/2\pi)\ln(\mu/\mu_0)$.

The question is whether any derivation path adds an extra factor of $\Nhelicity = 2$
on top of this baseline.

%──────────────────────────────────────────────────────────────────────────────
\section{Path A — Fermionic Dirac Trace}
\label{sec:pathA}
%──────────────────────────────────────────────────────────────────────────────

\subsection{Statement of Path A}

Route~R1 (\texttt{neff12\_derivations.tex} §R1 step 2) states: "The
complexification $\mathbb{C}\otimes$ provides a complex structure on $\mathbb{H}$,
splitting each quaternion field component into left and right Weyl helicity
sectors.  Thus $\Nhelicity = 2$."

\texttt{step2\_vacuum\_polarization.tex} §3.3 claims: "The standard
Passarino-Veltman reduction gives a factor $\Nhelicity = 2$ relative to a
scalar loop without polarization sum."

\subsection{Analysis}

The factor of 2 in a fermion loop (relative to a complex scalar loop)
comes from the Dirac trace:
\begin{equation}
  \mathrm{Tr}\!\left[\slashed{p}\,\gamma^\mu\,(\slashed{p}+\slashed{k})\,\gamma^\nu\right]
  = 4\!\left[p^\mu(p+k)^\nu + p^\nu(p+k)^\mu - g^{\mu\nu}p\cdot(p+k)\right].
  \label{eq:Dirac_trace}
\end{equation}
The factor of 4 from the Dirac trace (versus 1 for a complex scalar) accounts
for the spin-1/2 degrees of freedom.  For a Weyl fermion the trace gives a
factor of 2 instead of 4.

Passarino-Veltman reduction is a tensor-reduction technique that expresses
loop integrals with Lorentz tensor structure in terms of scalar integrals
$B_0, B_1, \ldots$.  It does \emph{not} introduce helicity multiplicities.
Helicity multiplicities arise \emph{before} PV reduction, from the trace over
spin indices in the numerator of the loop integrand.

\subsection{Does Path A apply to $S_{\mathrm{kin}}[\Theta]$?}

The action \eqref{eq:S_eff} is a \emph{scalar} kinetic term.  There is no
Dirac matrix $\gamma^\mu$ in the numerator of the $\Phi_k$ propagator.
The loop integrand is
\begin{align}
  \Pi^{\mu\nu}(k) &\propto \int \frac{d^4p}{(2\pi)^4}
  \frac{(2p+k)^\mu(2p+k)^\nu}{[p^2-m^2][(p+k)^2-m^2]},
  \label{eq:scalar_loop}
\end{align}
and there is no trace over spinor indices.  The resulting scalar-QED loop
coefficient is exactly the baseline \eqref{eq:baseline}; no additional
factor of 2 appears.

\begin{proposition}[Path A does not apply]
\label{prop:pathA}
Path A is inapplicable to $S_{\mathrm{kin}}[\Theta]$ as written.  The
minimal action is a scalar action with no Dirac fermion term.  Therefore
$\Nhelicity = 2$ cannot be derived from a Dirac trace argument for this action.
\end{proposition}

%──────────────────────────────────────────────────────────────────────────────
\section{Path B — Off-Diagonal $\mathrm{Mat}(2,\mathbb{C})$ Components}
\label{sec:pathB}
%──────────────────────────────────────────────────────────────────────────────

\subsection{Statement of Path B}

Route~R4 (\texttt{neff12\_derivations.tex} §R4) uses the isomorphism
$\mathbb{C}\otimes\mathbb{H} \cong M_2(\mathbb{C})$ and counts:
\begin{equation}
  \Neff = \underbrace{2}_{\text{off-diag}}
         \times
         \underbrace{3}_{N_{\mathrm{phases}}}
         \times
         \underbrace{2}_{\Nhelicity}
  = 12.
  \label{eq:R4}
\end{equation}
The "two Weyl helicity components from the $\mathbb{C}$ tensor factor"
are supposed to give $\Nhelicity = 2$.

\subsection{Analysis}

In $M_2(\mathbb{C})$, the two off-diagonal entries $\theta_{12}$ and
$\theta_{21}$ are charged under $U(1)_{\mathrm{EM}}$ with charges $+q$
and $-q$ respectively.  The two diagonal entries $\theta_{11}$ and
$\theta_{22}$ are neutral.

Under the minimal coupling $D_\mu = \partial_\mu + ieA_\mu$ acting on
the whole matrix $\Theta$, the kinetic term
\begin{equation}
  \mathrm{Tr}\!\left[(D_\mu\Theta)^\dagger(D^\mu\Theta)\right]
  \label{eq:Tr_kin}
\end{equation}
splits into contributions from all four components.  The charged
off-diagonal entries each contribute:
\begin{equation}
  \bigl|({\partial}_\mu + ieA_\mu)\theta_{12}\bigr|^2
  + \bigl|({\partial}_\mu - ieA_\mu)\theta_{21}\bigr|^2.
  \label{eq:offdiag}
\end{equation}
This is two independent complex scalars with charges $\pm e$.

\begin{proposition}[Path B gives $N_{\mathrm{charge}} = 2$, not $\Nhelicity = 2$]
\label{prop:pathB}
The factor of 2 from the two off-diagonal entries in Path~B accounts for
the charge-conjugate pair $(\theta_{12}, \theta_{21})$.
This is the particle/antiparticle doubling already incorporated in the
standard complex-scalar QED coefficient $e^2/(6\pi^2)$
(see \texttt{step2\_vacuum\_polarization.tex} §3.4).
It is $N_{\mathrm{charge}} = 2$, not $\Nhelicity = 2$.
Adding it as an independent factor on top of the standard complex-scalar
formula is double counting.
\end{proposition}

\begin{remark}
The "two Weyl helicity components from the $\mathbb{C}$ tensor factor"
in Route~R4 cannot be identified with left/right helicity states without
a separate kinetic term coupling $\Theta$ as a Weyl spinor (e.g.,
$\bar\Theta\, i\sigma^\mu D_\mu \Theta$ for a left-Weyl field).
The minimal scalar action \eqref{eq:S_kin} does not contain such a term.
\end{remark}

%──────────────────────────────────────────────────────────────────────────────
\section{Path C — Weyl Decomposition from $\mathbb{C}\otimes\mathbb{H}$}
\label{sec:pathC}
%──────────────────────────────────────────────────────────────────────────────

\subsection{Statement of Path C}

The complexification $\mathbb{C}\otimes\mathbb{H}$ provides a natural Weyl
decomposition because $\mathbb{H}$ contains the quaternionic anti-involution
$\Theta \mapsto -\mathbf{k}\Theta\mathbf{k}$ which (in the Mat$(2,\mathbb{C})$
representation) corresponds to charge conjugation.  One could argue that each
quaternionic $\Phi_k$ is a Weyl spinor with two helicity components $\phi_k^L$
and $\phi_k^R$, contributing a factor of $\Nhelicity = 2$.

\subsection{Analysis}

For this argument to give a loop multiplicity of 2, the fields $\phi_k^L$ and
$\phi_k^R$ must appear as \emph{independent} loop propagators in the Feynman
diagram for $\Pi^{\mu\nu}$.  This requires a kinetic term of the form:
\begin{equation}
  \mathcal{L}_{\mathrm{Weyl}}
  = i\phi_k^{L\dagger}\bar\sigma^\mu({\partial}_\mu + ieA_\mu)\phi_k^L
  + i\phi_k^{R\dagger}\sigma^\mu({\partial}_\mu + ieA_\mu)\phi_k^R.
  \label{eq:Weyl_kinetic}
\end{equation}
Such a term is \emph{not present} in $S_{\mathrm{kin}}[\Theta]$ as written.
Deriving \eqref{eq:Weyl_kinetic} from \eqref{eq:S_kin} would require:
\begin{enumerate}
  \item An identification of the Weyl components $\phi_k^{L/R}$ within $\Phi_k$
        as independent propagating fields (not just algebraic components).
  \item A proof that the kinetic operator of $\Phi_k$ in the winding background
        reproduces \eqref{eq:Weyl_kinetic}, not the scalar kinetic term
        $|({\partial}_\mu + ieA_\mu)\Phi_k|^2$.
\end{enumerate}

Neither of these steps is carried out in the current canonical derivation chain.

\begin{proposition}[Path C requires additional derivation]
\label{prop:pathC}
The Weyl-decomposition argument for $\Nhelicity = 2$ from $\mathbb{C}\otimes\mathbb{H}$
is physically motivated by the algebraic structure but is not derived from
$S_{\mathrm{kin}}[\Theta]$.  A separate fermionic effective action must be derived
from $S[\Theta]$ (e.g.\ via a Dirac operator from the biquaternionic structure)
before $\Nhelicity = 2$ can be upgraded to proved status.  This is a future
derivation task, not a current result.
\end{proposition}

%──────────────────────────────────────────────────────────────────────────────
\section{Loop-Safe Verdict}
\label{sec:verdict}
%──────────────────────────────────────────────────────────────────────────────

\begin{theorem}[$\Nhelicity = 1$ is the loop-safe value from $S_{\mathrm{kin}}[\Theta]$]
\label{thm:nhelicity}
For the minimal UBT kinetic action $S_{\mathrm{kin}}[\Theta]$, the one-loop
vacuum polarisation of the photon is computed from three copies of scalar QED
for the complex fields $\Phi_k$ ($k = 1, 2, 3$).  No independent helicity
multiplicity factor arises from this action.  The loop-safe value is:
\begin{equation}
  \boxed{
  \Nhelicity^{\mathrm{loop}} = 1
  \quad\Rightarrow\quad
  \Neff(\mathrm{loop}) = N_{\mathrm{phases}} \times \Nhelicity^{\mathrm{loop}}
  = 3 \times 1 = 3
  \quad\Rightarrow\quad
  B_0 = \frac{2\pi}{3}\times 3 = 2\pi.
  }
  \label{eq:verdict}
\end{equation}
\end{theorem}

\begin{proof}
Path A: inapplicable — no Dirac trace (Proposition~\ref{prop:pathA}).\\
Path B: gives $\Ncharge = 2$, already included in the standard complex-scalar
coefficient; adding it again is double counting (Proposition~\ref{prop:pathB}).\\
Path C: requires a separate fermionic effective action not present in
$S_{\mathrm{kin}}[\Theta]$ (Proposition~\ref{prop:pathC}).

In the absence of a valid derivation of $\Nhelicity = 2$ from
$S_{\mathrm{kin}}[\Theta]$, the loop calculation produces exactly
$3 \times B_0^{(1)} = 3 \times (2\pi/3) = 2\pi$.
\end{proof}

%──────────────────────────────────────────────────────────────────────────────
\section{Update to step2\_AUDIT.tex §Step B Verdict}
\label{sec:update_audit}
%──────────────────────────────────────────────────────────────────────────────

The verdict in \texttt{step2\_AUDIT.tex} §Step~B should be read as:

\begin{itemize}
  \item $\Nhelicity = 2$ is \textbf{not proved} as an independent loop multiplier
        from $S_{\mathrm{kin}}[\Theta]$.
  \item Three candidate derivation paths (A, B, C) have been examined.  None
        produces $\Nhelicity = 2$ as an independent multiplier:
        \begin{itemize}
          \item Path A (Dirac trace): requires fermionic kinetic term not in $S_{\mathrm{kin}}$.
          \item Path B (matrix structure): produces $\Ncharge = 2$, not $\Nhelicity = 2$.
          \item Path C (Weyl decomp.): requires derivation of fermionic effective action.
        \end{itemize}
  \item The loop-safe value $\Nhelicity = 1$ gives $\Neff = 3$, $B_0 = 2\pi$.
  \item Status: \textbf{OPEN / [MC]} — not closed without a derived fermionic
        action from $S[\Theta]$.
  \item Dedicated derivation audit:
        \texttt{research\_tracks/quantum\_ubt/fermionic\_action\_derivation.tex}.
\end{itemize}

\begin{statusbox}
Výsledek: $\Nhelicity = 1$ (loop-safe from $S_{\mathrm{kin}}[\Theta]$);
$\Nhelicity = 2$ is unproved as an independent multiplier.\\
Klasifikace: [MC] OPEN for $\Nhelicity = 2$; [L1] proved for $\Nhelicity = 1$
  from scalar action.\\
Dopad na alpha track: $\Neff(\mathrm{loop}) = 3$, $B_0 = 2\pi$, $n^*\approx 10.54$.
Upgrading $\Nhelicity$ to 2 requires a new derivation of the fermionic
sector from $S[\Theta]$ (future task; see
\texttt{research\_tracks/quantum\_ubt/fermionic\_action\_derivation.tex}).
\end{statusbox}

\end{document}
